\subsection{The Recurrent State-Space Model and the Dreamer Line}
\label{section:fc_rssm}

The previous section described the modular Vision-Memory-Controller architecture of \citet{HaSchmidhuber2018}, trained in three separate stages. The Recurrent State-Space Model collapses that sequence into a single jointly trained network and becomes the workhorse architecture of modern world-model reinforcement learning. The line carries through PlaNet, DreamerV1, DreamerV2, and DreamerV3 and is treated here in version-agnostic notation; the differences across versions are summarized at the end of the subsection. The pivot at the end is to the question of what loss the model should be trained against, which the next section picks up.

\subsubsection{The Recurrent State-Space Model}
\label{section:fc_rssm_transition}

Let $o_t$ denote the observation at time $t$, $a_t$ the action, and $r_t$ the reward. The Recurrent State-Space Model factors the latent state into a deterministic recurrent component $h_t$ and a stochastic latent $s_t$, coupled by
\[ h_t = f_\theta(h_{t-1}, s_{t-1}, a_{t-1}), \qquad s_t \sim p_\theta(s_t \mid h_t). \]
The deterministic path $f_\theta$ is a recurrent cell (in practice a Gated Recurrent Unit); the prior $p_\theta(s_t \mid h_t)$ is either a diagonal Gaussian or a product of categoricals, the choice a design parameter rather than a structural commitment. A posterior network $q_\theta(s_t \mid h_t, o_t)$ encodes the current observation into a refined latent for inference at training time. The joint loss for a length-$T$ trajectory has three named terms,
\[ \mathcal{L}_\theta = \underbrace{-\sum_t \log p_\theta(o_t \mid h_t, s_t)}_{\text{reconstruction}} \;-\; \underbrace{\sum_t \log p_\theta(r_t \mid h_t, s_t)}_{\text{reward log-likelihood}} \;+\; \underbrace{\sum_t \mathrm{KL}\!\left[ q_\theta(s_t \mid h_t, o_t) \,\Vert\, p_\theta(s_t \mid h_t) \right]}_{\text{prior-posterior consistency}}. \]
Optionally a latent-overshooting term enforces multi-step consistency in latent space without decoding. The point of the architecture is to separate deterministic memory from stochastic uncertainty so that the model forecasts a distribution over futures while still propagating an information bottleneck through time. The model class is a flexible neural object, but the algorithmic outer loop is recognizably the Dyna template of \S\ref{section:fc_dyna_q}, with $\widehat P$ and $\widehat r$ replaced by jointly trained networks.

\subsubsection{Actor-critic in imagination}
\label{section:fc_rssm_imagination}

The policy is trained inside the learned world model rather than from environment interaction directly. An actor $\pi_\phi(a_t \mid h_t, s_t)$ and a critic $V_\psi(h_t, s_t)$ are updated on imagined trajectories rolled forward through the latent dynamics, starting from real states drawn from the agent's replay buffer. The imagined return objective is
\[ J(\phi) = \mathbb{E}_{\pi_\phi, p_\theta}\!\left[ \sum_{\tau = t}^{t+H} \gamma^{\tau - t} \hat r_\tau + \gamma^{H+1} V_\psi(h_{t+H+1}, s_{t+H+1}) \right], \]
where the expectation is taken over actions sampled from the actor and latent states evolved by the world model, $H$ is an imagination horizon, $\hat r_\tau$ is the reward predicted by the trained reward head, and the bootstrap term $V_\psi$ supplies a value estimate past the imagination horizon. Because the world-model dynamics are differentiable and the latent samples are reparameterized, $\nabla_\phi J$ is obtained by backpropagating value gradients through the unrolled latent trajectory. The critic is fit by minimizing a squared error against $\lambda$-return targets in the same imagined trajectories.

This is the deep-learning instantiation of the inner planning loop of Dyna-Q, with the model now a flexible neural object and the planner an analytic policy-gradient operator rather than a discrete Q-update. Convergence guarantees are not available in closed form for this setting; the nearest analogue is the Lipschitz contraction bound for model-based value error of \citet{Asadi2018Lipschitz}, discussed in \S\ref{section:fc_synthesis_fails}, which states that value error grows geometrically along the rollout at rate $K_F$ (the Lipschitz constant of the learned dynamics) and remains finite when $\gamma K_F < 1$.\footnote{The four canonical instantiations of the architecture cover the major design choices. PlaNet \citep{HafnerPlaNet2019} introduces the RSSM with cross-entropy planning. DreamerV1 \citep{HafnerDreamer2020} replaces planning with actor-critic in imagination. DreamerV2 \citep{HafnerDreamerV2_2021} switches the Gaussian latent for categoricals with a Kullback-Leibler balancing trick, surpassing top model-free agents at the Atari 200M budget. DreamerV3 \citep{HafnerDreamerV3_2025} fixes one hyperparameter configuration across more than one hundred and fifty tasks via symlog reward normalization and free-bits Kullback-Leibler, including the Minecraft diamond from scratch.}

The line trains the world model against likelihood-style losses (reconstruction, reward log-likelihood, prior-posterior consistency), which rewards the model for matching observed distributions irrespective of how those distributions enter the value calculation. Read in economics terms, the recurrent state-space model is a high-dimensional perceived law of motion, fit by likelihood rather than by hand-chosen moments, with a planner that operates entirely on the learned latent. The next section asks whether this is the right objective. The argument, developed by \citet{FarahmandVAML2017} and \citet{GrimmVALEQ2020} and instantiated empirically in MuZero, is that the model should be trained against the value functional directly rather than against pixel and reward likelihoods. The Dreamer line achieves its results despite training to the wrong target; the value-aware line asks how much further the field could go if the loss were aligned with the downstream decision.
